\section{Assumptions, conditions, and unavailable authority}
\label{sec:assumptions}

The equations are conditional on the following explicit statements. They are
not permissions to replace missing information with defaults.

\begin{description}[style=nextline]
\item[\Assump{001}: upstream sample authority.]
  The supplied sample axis, time, eligibility, original/synthesized state,
  elevation/airmass meaning, and acquisition-column binding are correct under
  approved upstream authority. SCI-CAL neither chooses nor rederives them.

\item[\Assump{002}: ordinary measured channel.]
  The admitted signal is the ordinary \code{xs} stream. A distinct measured
  stream is a distinct physical observable and does not inherit this factor.

\item[\Assump{003}: selected immutable APT.]
  TolProj has selected and bound one immutable measured APT to the target
  observation, and the supplied target-to-source associations are admitted
  TolAPT/TolProj results. SCI-CAL does not choose the artifact or matcher.

\item[\Assump{004}: explicit affine convention.]
  Any baseline or offset operation preceding the multiplier is named, typed,
  and already authorized. ``No baseline subtraction'' is an explicit
  not-applicable state, not an inferred zero.

\item[\Assump{005}: coefficient direction and reference plane.]
  The selected \code{flxscale} is finite, nonzero, oriented as
  \(\mjybeam/U_x\), and refers to the top-of-atmosphere calibration plane. Its
  lineage states the disposition of any TolProj pointing-derived correction.

\item[\Assump{006}: atmosphere inputs.]
  \(\tau_{225}\) is zenith optical depth; $X$ is the full eligible-sample
  airmass; the target pivot is $X_{\rm ref}=0$; and every value carries
  declared time and model support. SCI-CAL does not infer $X$ from elevation
  or choose an opacity estimator.

\item[\Assump{007}: fixed conditional state.]
  Conditional covariance equations hold at fixed identities, factor values,
  atmosphere state, selection, and response state. Dependence of fitted state
  on the same signal requires cross-covariance or joint propagation.

\item[\Assump{008}: response ownership.]
  Beam/Beammap supplies the originating beam/template identity; MAP/FLT
  supplies the realized mapmaker/kernel/filter response. SCI-CAL records and
  conditions on those identities but does not certify their empirical
  fidelity.

\item[\Assump{009}: coherent segment.]
  A segment is declared before quality classification and is internally
  coherent in selected APT, atmosphere/operator identity, requested/effective
  mode, target unit, and response basis. Splitting a segment is a provenance
  event, not a sample-by-sample quality switch.

\item[\Assump{010}: passband limitation.]
  The selected TolTECA v1 passband set is a modeled array-named reference with
  the recorded unknowns in Section~\ref{sec:definitions}. Any atmosphere
  result using it is conditional on the declared, presently unestablished
  response weighting and carries unavailable passband uncertainty.

\item[\Assump{011}: unresolved numeric operator authority.]
  The packet does not contain the exact nodes, node ordinate orientation,
  content digest, or numeric line-of-sight support of
  \code{am12_fixed_djf25_piecewise_linear_los_tau_v1}. Consequently this draft
  defines the strongest structural operator contract but does not authorize a
  numeric atmosphere evaluation, calibrated numeric output, or numeric
  \code{science-qualification-eligible} disposition until the owner supplies
  one content-bound operator record answering the question in the draft
  decision register. Resolving that question can enable structural evaluation
  only; it does not establish an achieved \code{science-qualified} or
  \code{calibrated-science} claim.
\end{description}
